\section{Discussion}
\label{sec:discussion}

The results support a disciplined reinterpretation of DeFi security evidence. The incident layer says that DeFi losses are numerous, recurrent, and diverse. The protocol layer says that pure on-chain exploitability among major protocols is much narrower. The replay layer says that even historically important incidents do not all translate cleanly into executable proofs without substantial engineering and runtime context. The contribution of the workflow is therefore epistemic as much as empirical: it forces each claim to be attached to the evidence layer that actually supports it.

Three mechanism families dominate the current artifact set. Oracle manipulation remains the clearest bridge between protocol design assumptions and replayable exploitation, especially when external price-reporting paths or thin-liquidity dependencies interact with permissive borrowing logic. Flash-loan-assisted logic abuse remains important because it converts otherwise awkward state transitions into atomic capital-efficient strategies. Reentrancy remains relevant not because every protocol is equally vulnerable, but because callback semantics and external control flow still create sharp discontinuities in otherwise mature systems. These families are consistent with prior DeFi security literature \citep{qin2021attacking,zhou2023sok}, but the repository adds concrete workflow scaffolding around them.

The audit and verification layers also suggest a practical distinction between \emph{security debt} and \emph{replayable exploitability}. A mature protocol may still accumulate dozens of medium-severity hardening issues or architectural concerns without exposing a high-confidence pure on-chain exploit path. This is visible in the five deep audits, where 87 manually triaged findings coexist with a protocol-verification result in which only one out of 24 major protocols is classified as confirmed exploitable. Research papers that ignore this distinction risk confusing engineering debt with imminent exploitability. Conversely, a workflow that only studies confirmed exploits may ignore the long tail of design weaknesses that shape future incidents.

The main limitation of the present paper is that the replay layer is still incomplete. The three principal replay cases have concrete logic encoded, but archive RPC validation has not been completed in the current environment. The protocol-verification layer also contains conditional cases whose exploitability depends on downstream integrations or edge-case market conditions. These limits do not invalidate the paper's argument; they define its scope. The paper is strongest as a methodology and intermediate empirical study, not yet as a definitive census of fully reproduced DeFi exploits.

That limitation points directly to the next phase of work. The repository should expand the replay suite from three implemented cases to at least five or eight archive-validated cases spanning oracle manipulation, reentrancy, logic abuse, and governance-assisted paths. It should also improve linkage between incident cards, protocol studies, and replay prerequisites so that each paper claim can be traced to a stable artifact. The refactored writing system created here is intended to make that next step cheaper, because it separates experiment generation, paper drafting, and revision rather than collapsing everything into a single ad hoc document.
